% !TeX root = 20261020_context_design.tex
% !LW recipe=lablog XeLaTeX
\input{_support/style.tex}
\begin{document}
\reporttitle{2026-10-20}{1,000개 버퍼에 맞춘 소수 신규 공격 표본 수 비교 설계}
\recordnote{2026-10-20 사용자 피드백을 반영한 개정 설계. 같은 날 사용자가 100개 조건 추가와 실행을 요청하여 연구 버전 v7의 EXP84로 시작했다. 진행 상태·관찰 결과는 \path{reports/20261020_exp84.tex}에 기록한다. 이번 개정은 신규 공격당 공급 수 자체를 조정하는 것이 핵심이다. 앞선 편입 상한·클래스별 보호량·ICD pilot 계획을 대체하며, 이미 수행한 EXP83의 결과와 설정은 변경하지 않는다.}

\section{연구 질문과 수정 범위}
새 공격의 라벨 사례가 소량 확보되었을 때, \textbf{XGB는 누적 사례로 재학습하고 TabPFN은 사전학습 가중치를 고정한 채 context를 갱신하면 각각 어느 정도의 탐지 품질과 갱신 비용을 보이는가}를 비교한다. TabPFN에서는 모든 과거 사례를 누적하는 경우와 1,000개 버퍼를 유지하는 경우를 함께 본다. 낮은 비용이나 기존 공격 성능 유지가 자동으로 보장된다는 주장을 전제로 하지 않는다.

이전 설계는 10--400개라는 라벨 공급량을 유지하면서 실제 편입량만 제한했다. 예를 들어 400개를 확보해도 일부만 넣도록 하여, 사용자가 확인하려는 ``실제로 몇 개의 신규 사례가 필요한가''와 다른 선택·압축 문제를 추가했다. 이번에는 \textbf{공급한 신규 사례를 모두 사용}하고, 공급 수 자체를 5--50개 중심으로 낮추되 사용자 요청에 따라 100개도 비교한다. 정상 500개·공격 25개 보호, 단계당 편입 상한, 신규 대표점 선택, 동일 버퍼 XGB 대조, ICD pilot은 이번 비교에서 제외한다. 버퍼 크기로 결정되는 교체율 그래프도 삭제하고 설정 숫자는 표로만 제시한다.

기존 신규 공격당 100개는 설명용으로 골랐던 조건이며 최적값·통계적 중앙값·현실의 하루 공격량이라는 근거가 없다. 이전 결과의 한 예로만 설명하고, 필요 표본 수는 곡선 전체로 판단한다.

\section{조정할 숫자와 모든 신규 사례의 투입}
초기 사례와 버퍼 용량은 모두 1,000개로 유지하고 신규 공격 \emph{클래스당} 사례 수를
\[
 n\in\{5,10,20,30,50,100\}
\]
으로 바꾼다. 한 단계에 새 클래스가 $k_t$개면 모든 방법에 같은 $k_tn$개를 제공한다. 고정 버퍼는 기존 사례에서 정확히 $k_tn$개를 삭제하고 새 사례를 전부 넣는다. 실제 편입 수를 줄이는 별도 상한이나 신규 클래스당 제한은 없다.

\begin{table}[H]\centering\small
\caption{새 비교의 설정값. 교체율은 설계로 결정되는 값이며 측정 결과가 아니다. 누적 수는 모든 공격 도입 후 XGB 학습/PFN context에 사용하는 총 사례 수다.}
\begin{tabular}{rrrrr}
\toprule
신규 공격당 $n$ & CIC 단계당 교체 & ToN 최대 단계 교체 & CIC 최종 누적 & ToN 최종 누적\\
\midrule
5  & 5 (0.5\%)  & 10 (1.0\%)  & 1,020 & 1,035\\
10 & 10 (1.0\%) & 20 (2.0\%)  & 1,040 & 1,070\\
20 & 20 (2.0\%) & 40 (4.0\%)  & 1,080 & 1,140\\
30 & 30 (3.0\%) & 60 (6.0\%)  & 1,120 & 1,210\\
50 & 50 (5.0\%) & 100 (10.0\%) & 1,200 & 1,350\\
100 & 100 (10.0\%) & 200 (20.0\%) & 1,400 & 1,700\\
\bottomrule
\end{tabular}
\end{table}
CIC2018은 초기 이후 총 4개 공격을 한 클래스씩, ToN-IoT는 총 7개 공격을 단계당 최대 두 클래스씩 도입한다. 따라서 최종 누적 사례 수는 각각 $1{,}000+4n$, $1{,}000+7n$이고, 고정 버퍼는 전 단계에서 1,000개다.

5--50개 범위는 ToN의 두 클래스 동시 도입에서도 한 번에 버퍼의 10\% 이하를 교체한다. 사용자 요청으로 추가한 100개 조건에서는 최대 200개(20\%)를 교체한다. 5개는 기존 최저 조건 10개보다 더 적은 사례를 확인하고, 10·20·50·100개는 기존 결과와 연결하며 30개는 그 사이를 보완한다. 이 범위가 성능을 보장하거나 최적이라는 뜻은 아니다. ToN의 $n=50/100$은 여러 단계를 합하면 각각 350/700회의 사례 삭제가 발생하므로, 단계당 교체가 작아도 기존 성능 유지 여부는 측정해야 한다.

이전 EXP83의 $n=400$에서는 CIC 단계당 40\%, ToN 최대 80\%를 교체했다. FIFO의 최종 정상/초기 공격 보유 수는 두 데이터 모두 1/2개였고, 같은 $n=400$에서도 정책별 점수 차이가 컸다. 큰 교체량과 기존 사례 감소를 함께 관찰했지만 그것만을 단일 원인으로 확정하지 않는다. 이번에는 삭제 규칙을 그대로 두고 표본 수 범위를 낮춰 관찰한다.
\sourcepath{docs/lablog/data/exp83_budget_values.csv}
\sourcepath{configs/experiments/exp83_protocol.json}

\section{공통 설정·비교군·평가와 비용}
\subsection{데이터와 단계}
CIC2018과 ToN-IoT를 각각 독립적으로 평가한다. 정제된 각 벤치마크의 기존 train/test 분리를 유지하고, 같은 벤치마크의 train에서만 초기·신규 사례를 선택한다. 다른 데이터셋의 공격을 옮겨 넣거나 새로운 특징 행을 합성하지 않는다. 설계한 부분은 공격 라벨을 공개하는 단계와 그 수량이다. 원본 트래픽 전체를 시간순으로 재생하는 실험이나 라벨 확보 전 미지 공격 탐지 실험은 아니다.

초기는 정상 750개와 기존 공격 두 클래스 각각 125개로 총 1,000개다. 기존 초기 행 ID를 그대로 사용한다. 아래 표는 \textbf{공격 도입 단계순}이며, 각 칸은 그 단계에 처음 제공하는 학습/context 사례를 뜻한다.
\begin{table}[H]\centering\small
\begin{tabular}{cll}
\toprule
단계 & CIC2018 & ToN-IoT\\\midrule
1 & 정상 750, Brute force 125, DoS 125 & 정상 750, DoS 125, Scanning 125\\
2 & DDoS $n$ & DDoS $n$, Injection $n$\\
3 & Web attacks $n$ & Password $n$\\
4 & Infiltration $n$ & XSS $n$\\
5 & Bot $n$ & Backdoor $n$, Ransomware $n$\\
6 & 해당 없음 & MitM $n$\\\bottomrule
\end{tabular}
\end{table}
seed는 43 한 번이며 반복 seed는 수행하지 않는다. EXP78과 같은 클래스별 중첩 표본 추출을 사용하여 10·20·50개 조건의 행 ID를 기존 실행과 일치시키고 5·30개를 같은 순서에서 선택한다. 100개는 EXP77의 원래 계획과 일치시킨다. 모든 비교군은 각 단계에 정확히 같은 신규 행을 받는다. 별도 validation 라벨을 추가하지 않는다. 기존 test는 원본 행 ID가 train과 분리되어 있지만 동일 특징 벡터까지 모두 분리된 집합은 아니므로, 기존에 보고한 특징 중복 제외 진단도 별도로 보존한다.

\subsection{비교군과 버퍼 규칙}
\begin{table}[H]\centering\small
\begin{tabular}{p{.24\linewidth}p{.68\linewidth}}
\toprule
비교군 & 단계별 동작\\\midrule
XGB 누적 재학습 & 현재까지 받은 사례 전체로 매 단계 모델 재학습\\
PFN 누적 context & 같은 사례 전체를 context로 사용하여 전처리·캐시 준비\\
PFN FIFO & 가장 먼저 들어온 기존 사례부터 삭제 후 신규 전량 투입\\
PFN 무작위 & 기존 사례를 seed 고정 무작위 순서로 삭제 후 신규 전량 투입\\
PFN 유사도 & 신규 사례와 가까운 기존 사례부터 삭제 후 신규 전량 투입\\\bottomrule
\end{tabular}
\end{table}
세 버퍼 정책은 EXP83의 구현을 유지한다. FIFO는 사례 획득 순서이며 실제로 낡았거나 무효해졌다는 의미는 아니다. 유사도는 현재 \emph{기존 context}에 맞춘 표준화 공간에서 신규 사례까지의 최소 유클리드 거리다. 무작위 seed는 기존과 같은 $43+1009s$이며 $s$는 0부터 시작하는 단계 번호다.

기존 코드의 클래스별 최소 1개 유지 규칙만 그대로 둔다. 이 규칙은 해당 클래스의 마지막 사례를 삭제 대상에서 건너뛰는 것으로, 신규 투입을 거절하는 규칙이 아니다. 새 범위의 최대 투입은 200개이고 최대 기존 클래스 수는 9개이므로, 1,000개에서 필요한 수만큼 삭제 가능한 여유가 있다. 새 보호량·클래스 균등화는 추가하지 않는다. 삭제된 사례는 이후 다시 가져오지 않는다.

모든 PFN은 EXP83과 같은 v3 checkpoint, estimator 4개, CPU context cache, \path{fit_with_cache}, context subsampling 없음을 유지한다. 초기 IDS 가중치 fine-tuning과 후속 fine-tuning은 모두 없다. 데이터셋 2개 $\times$ 표본 수 6개 $\times$ 비교군 5개로 \textbf{60조건}이다. LoCalPFN·앙상블·DistPFN·ICD를 추가하지 않는다.

\subsection{표본 수에 따른 탐지 품질과 기존 성능 유지}
주 결과는 마지막 상태의 모델로 동일한 전체 test를 평가한다. CIC2018은 3,042,473개, ToN-IoT는 2,271,723개다. seed 43으로 섞은 동일 query 순서와 배치 8,192개를 사용한다. 단계별 예측을 합산해 최종 성능으로 표시하지 않는다.

성능 그림의 가로축은 신규 공격당 $n$이며 데이터셋별로 다음을 구분한다. (1) 최종 전체 Macro-F1, (2) 초기 이후 도입한 공격들의 클래스별 F1 평균과 클래스별 recall/F1, (3) 갱신 전후 기존 클래스 성능 변화다. 최종 클래스별 F1은 \emph{전체 test}의 혼동행렬에서 계산하여 다른 클래스의 false positive를 포함한다. 기존 정상·공격에 대한 최종 점수만으로 ``유지''라고 표현하지 않는다.

유지 진단은 test에서 클래스당 최대 2,000개를 seed 43으로 뽑아 모든 단계에서 같은 행 목록·같은 혼합 query 순서로 예측한다. 각 갱신 직전까지 알려진 클래스의 동일 행에서 직전/직후 모델을 비교한다. 초기 세 클래스에 대한 단계별 점수도 별도 표시한다. 이 패널에는 미래 클래스의 query 특징이 있을 수 있지만 정답은 측정에만 쓰며, 오프라인 진단을 시간순 운영 재생이나 미지 공격 탐지로 해석하지 않는다. 기존 행만의 평가에는 신규에서 기존으로 발생하는 오탐이 빠지므로, 도입된 전체 클래스에 대한 기존/신규 precision·F1과 정상 FPR도 함께 기록한다.

60조건 모두 같은 환경에서 실행해 비용을 비교한다. 겹치는 10·20·50·100개 과거 예측은 재현 대조에 쓸 수 있으나 과거 측정 시간과 새 시간을 섞어 속도 우위를 주장하지 않는다. 단일 seed의 탐색 결과이므로 작은 점수 차이를 통계적 우위나 동등성으로 단정하지 않고, 사전 품질 기준이 없는 상태에서 ``최소 몇 개면 충분''이라고 확정하지 않는다.

\subsection{갱신 비용}
주 비용은 \textbf{초기 단계를 제외한 후속 갱신 시간의 합계}다. XGB는 사례 누적·모델 구성·전처리·재학습, PFN은 context 누적/삭제·모델 구성·전처리·표현/캐시 준비를 포함한다. GPU 동기화 후 같은 장치에서 순차 측정한다. 초기 준비, test 추론, checkpoint 저장, 검증·보고서 생성 시간은 분리해 기록하고 주 갱신 그래프에는 합산하지 않는다.

고정 버퍼에서도 바뀐 행만 부분 계산하는 구현은 아니다. 1,000개 context 전체를 다시 처리할 수 있으므로, 적은 신규 표본만으로 갱신 시간이 줄어든다고 미리 결론 내리지 않는다. 성능 곡선과 함께 후속 갱신 시간 대 $n$ 곡선을 제시한다. 버퍼 크기·교체율은 설정표, 실제 클래스별 보유 수는 보조 표로 남기며 결정된 버퍼 크기를 재확인하는 그래프는 만들지 않는다.

\section{ICD 참고 근거와 이번 설계에서의 역할}
\emph{In-Context Data Distillation with TabPFN}의 참고 근거는 공개 코드가 아니라 \textbf{원문 §3의 손실식, Figure 2의 gradient 흐름, §4의 최적화 설정}이다. 가중치 $\theta$를 고정하고 context 특징 $X_{\mathrm{dist}}$를 학습한다는 방법을 문헌으로 참고할 수 있다. 이번 확인에서 저자 공개 ICD 구현은 확인하지 못했으며, 저자 코드를 사용했거나 원 결과를 그대로 재현했다는 의미로 서술하지 않는다.

저장소의 \path{scripts/v7/exp75/context_distillation.py}는 이 아이디어를 v3에 맞춘 \textbf{자체 구현}이다. 실제로 backbone을 고정하고 context tensor만 Adam으로 갱신하며, 클래스 균형 query 추출과 validation NLL 선택을 사용한다. 원문의 validation AUC 선택과 다르고 TabPFN 버전도 다르므로 공식 구현 또는 충실한 원 실험 재현으로 부르지 않는다.

\paragraph{Discussion.}
ICD는 큰 학습 집합을 작은 context로 표현하는 접근을 제시하지만, 실제 신규 라벨을 몇 개만 받아도 분류가 가능한지를 비교하는 이번 질문과는 다르다. context 최적화에는 역전파·선택 비용도 발생한다. 따라서 이번에는 문헌상 대안으로만 정리하고 실험군·일정에 넣지 않는다. 향후 압축 자체를 연구할 때 원문 기반 자체 구현임을 밝히고 별도 대조를 설계해야 한다.
\href{https://arxiv.org/html/2402.06971}{Ma 등, In-Context Data Distillation with TabPFN, arXiv:2402.06971 (2024), §§3--4}

\section{실행 전 검증과 발표 연결}
표본 추출과 삭제 함수는 기존 소스의 해당 함수를 그대로 읽어 실제 train 행 ID로 점검한다. 새 범위의 모든 정책에서 신규 전량 편입, 1,000개 유지, 클래스별 최소 1개, 삭제 사례 재사용 없음, 같은 신규 행, 10·20·50·100개 기존 계획 일치 여부를 확인한다. 이는 \textbf{행 구성 검증}이며 모델 실행이나 성능 검증이 아니다. 결과는 아래 JSON에 기록한다. EXP84에서는 가중치 불변·예측·비용까지 검증하고 v7 실행 등록 및 결과 보고를 마감한다.

발표에서는 운영 중 소수 신규 라벨 확보 문제, 사전학습 가중치와 외부 context의 관계, 두 벤치마크별 독립 train/test 및 단계별 사례 선택 과정을 순서대로 보여준다. ``현재 데이터셋'' 대신 ``현재까지 제공한 라벨 사례''를 사용한다. 완료된 10--400개 결과와 이번 5--100개 비교를 구별한다. 사용자는 결과 확인 후 슬라이드를 수정하기로 했으므로 EXP84 실행으로 발표자료를 변경하지 않는다. \path{docs/latex/main.tex}는 이번 결과를 채택하지 않고 별도 요청한 참고문헌 보완만 수행한다.
\sourcepath{configs/experiments/exp84_protocol.json}
\sourcepath{configs/experiments/proposals/20261020_context_admission.json}
\sourcepath{docs/lablog/evidence/20261020/context_admission_feasibility.json}
\sourcepath{scripts/common/build_context_admission_design.py}
\end{document}
